% SPDX-FileCopyrightText: 2026 Will Cook
% SPDX-License-Identifier: CC-BY-4.0
\section{Keeping a useful practice small}\label{sec:practice}
\papersectiontarget{exposition-practice}

Use the two-page \emph{Writing a Good Mathematical Paper}~\cite{compact} as a
working guide; use this companion when a particular instruction needs an
example or a limit. The repository's writing skill~\cite{skill} incorporates
the literature-reading pass, while its records retain the detailed cases and
exceptions. A writer need not read every past return before explaining a proof.
The next revision should resolve a specific difficulty, not demonstrate that
every rule has been mentioned.

\subsection{Revise in an order that preserves the argument}\label{sec:revision-order}

Begin with the full statement, proof and closest prior results. Write a brief
account of the contribution, the decisive inference and the exact boundary.
Check this account against the abstract and introduction. When a condition is
only sufficient, a witness is shared, or a parameter is fixed, preserve that
qualification in every compressed occurrence. Agreement among several summaries
is not enough if all have inherited the same overstatement.

Next choose the passage where the intended reader must supply the most
consequential missing explanation. Ask a concrete question: why this parameter,
why this normalisation, why this error threshold, or why does this limit give
the asserted object? Read a relevant local argument in an original source,
then find the mathematical answer in the manuscript. Write the sentence that
connects that answer to the next step. If the answer is absent, record the
mathematical issue rather than concealing it with a connective.

Review the revised paragraph with its predecessor and successor. Definitions
must arrive before they are used, and the paragraph's conclusion must supply
what the next one needs. Read displayed formulas as parts of grammatical
sentences, checking punctuation and the referent of each symbol. Conrad's
examples illustrate how unspecified variables and loose quantification can
make ordinary-looking prose say the wrong thing~\cite[\S\S1--2]{conrad}.
A sentence-level pass is valuable after the dependency structure is sound;
it cannot substitute for that earlier work.

Finally, remove explanation that repeats an already intelligible operation.
Keep a second account when it serves a different purpose, such as showing how
to apply a formula after deriving it. Stop when another sentence would resolve
no specific ambiguity, justify no missing transition and improve no useful
route through the paper. This gives brevity a reader-facing purpose without
imposing a page target.

\subsection{Keep the short paper and long record useful separately}\label{sec:short-long}

Before moving a passage, identify what a reader of each document must still be
able to do. The short paper should support its principal result through an
intelligible proof or a clearly identified proof sketch with a precise full-proof
destination. A long record should preserve the omitted derivation, its assumptions
and its relation to the main argument. A link labelled ``details'' is insufficient
when its destination proves a different statement or omits the difficult step.

Compare the shared assertions in both directions. A repaired endpoint in the
short paper must reach the long proof; an assumption exposed in the long proof
must reach the short statement. After a move, follow the link to the actual
passage and check that its notation can be translated. Retain substantial
alternatives and counterexamples where they explain the scope or a failed
method. Material that serves only the revision history can remain in the
editorial record rather than interrupting the mathematical argument.

\subsection{Put verification details where they can be checked}

The mathematical argument should state its hypotheses and explain the
inference. An exact computation may be part of that inference; give its
mathematical input, output and finite scope at the point of use. Put the
software version, source identifiers, replay commands and review history in
one verification and reproducibility section or appendix. Refer there from
a theorem when its evidence class matters. This keeps proof status available
without making a reader decode repository machinery between proof steps.

Give a single account of a shared limitation, then repeat it only when a
different claim would otherwise appear unconditional. A conditional premise
belongs beside the statement it conditions; an editorial disclaimer does
not need to recur after every equation. For an omitted proof, name its
mathematical destination by section or theorem and cite the original source
where an external result is used. A reference to a file or a ``long record''
alone does not tell a reader which argument to inspect.

\subsection{How the practice developed}

The recorded process began with reviewers proposing changes to particular
papers: move the result forward, explain a construction, shorten a repetitive
section or repair an unsupported implication. Successive returns did not form
a single consistent programme. In the earlier \#243 reviews, a cubic result
and a bounded negative result were each proposed as the lead at different
stages. Comparing these requests with the manuscripts yielded a narrower
instruction: choose the result that gives the current paper its most coherent
argument. The historical request remains evidence about that draft, not a
permanent ranking of result types.

Each returned manuscript then met a separate integration decision. In the
later second round, only one of nine candidates was accepted, although the
returns had passed transport checks. In the third round, integration included
restoring seventeen named-input remarks removed during compression. These
events exposed different questions: did the proposed files arrive intact,
did the revision preserve the mathematics and credit, and did it explain the
argument? A pass on the first question could not answer the others. Comparing
short papers with their longer records also recovered useful material, as the
\#251 buffer-coordinate case illustrates.

The recent rounds supplied more explicit literary models, before-and-after
passages and proposed general lessons. We compared these with the selected
native source, reopened the cited passages and retained repairs alongside
acceptances. The finite-remainder explanation survived this review; a claimed
Knuth locator needed correction; a broad author--reader analogy did not by
itself justify a specific rule about estimates. The resulting instructions
therefore combine mathematical checks with reading practice, while the
records preserve the disagreements and limits that compression would hide.

This short/long pair is the next distillation of that sequence. The compact
guide gives the writer actions; this companion explains why those actions are
useful and when they are insufficient. A later packet can carry both at a
fixed version, solicit another local revision and repeat the review. The
systems paper~\cite{systems} describes the repository mechanisms supporting
the process. Here the object of review is the editorial decision itself.

\subsection{From one case to a bounded amendment}

Most accepted changes should become examples of existing guidance. Add a rule
only when the case exposes a reusable distinction that the present instruction
misses. State the circumstance in which it applies, the action to take, the
evidence motivating it and a case in which it would fail. A page-one request
can become a reminder to remove unnecessary delay without becoming a page-one
quota. A failed proof can become an explanation of a missing hypothesis without
becoming a claim that every alternative approach is impossible.

The public records separate recent proposals, historical requests, source
identities and the manuscript decisions used to assess them. They retain
unmatched and pending cases explicitly. A return may describe an edit in
summary rather than supply the exact accepted words; this should be recorded
as a correspondence still requiring inspection, not counted as either a
rejected edit or a verbatim acceptance. Source passages are reviewed for their
stated writing use. Inclusion in the source list does not assert a complete
mathematical review of every cited paper.

The record of a failure should be equally specific. Preserve the attempted
implication, its assumptions, the witness to failure and the information
missing from the argument. Failure of one estimate does not refute the
theorem. Failure to run a check says that the check is unrun. An unavailable
source says something about the review's evidence, not about the truth of the
source's claim. These distinctions prevent a later writer from giving a
historical obstacle more force than it had.

\subsection{What the next packet must contain}

A frozen review packet should include the exact skill and guide it asks the
recipient to follow, the compact paper, the relevant lessons, the source
identities and this companion paper. A manifest records their paths and byte digests. The packet
must also carry the primary originals needed for its selected specimens,
with versions and reading-status declarations kept separate. A title in a
bibliography or a list of annex files is not a substitute for a source that
the recipient can actually inspect.

The recipient can then compare those papers with the supplied manuscript:
how does each introduce its object, make a hypothesis visible, motivate an
estimate or credit an input? The return should identify the passage read,
the feature adapted, the proposed local change and its limit. The repository
agent reviews the proposal against the current mathematics, integrates an
accepted change and amends the guidance only if the evidence warrants it.
The next packet can consume that resulting version. This is a sequence of
explicit document revisions, not training of model weights.

Freezing a version matters because guidance can change during a review.
An earlier packet cannot have applied a later correction to Knuth's locator,
and a later integration should not attribute that correction to the original
return. Keeping both versions makes the difference inspectable. The same
principle applies to mathematical status: a later accepted ordinary proof
does not turn an earlier unassessed proposal into an already checked result.

\subsection{Read the page that will be read}

The final pass uses the repository's native sources and rendering tools.
After changing a figure, inspect it at the size in which it appears. Follow a
reference to the intended theorem or long-form argument, rather than merely
checking that a target file exists. After reflow, inspect the affected pages
and their neighbours: a useful explanation can be separated from its figure,
or a small heading can be stranded above a page break. Bibliographic labels,
captions, interval endpoints and cross-document links belong to this review.

Read the result along three routes. First, scan the title, abstract,
introduction and main statements: recover the contribution and its boundary
without importing private knowledge. Second, reconstruct the proof: identify
where each hypothesis is used and explain the difficult inference in your own
words. Third, try to use a result: check its assumptions on an example, locate
an excluded case and find the cited input needed for an application.

For the examples here, these tasks have specific answers. In
Section~\ref{sec:certificate}, the certificate places a value in $(1/2,1)$;
its two inequalities come from the two endpoint gaps; the value $D=3/4$ with
the loose enclosure $[1/2,1]$ shows why a failed test is inconclusive. A reader
who can repeat ``compare both endpoints'' but cannot explain that failure has
not yet recovered how to use the certificate. This identifies a passage to
revisit, not a numerical score for understanding.

When an independent reader is available, ask for the exact place where their
reconstruction stopped and what they believed had been established there.
An author's own cold-start pass is still useful, but should be reported as
self-review. No independent reader is implied by the procedure described here.

The checks answer limited questions. A successful build establishes that the
source rendered under the recorded toolchain. A passage comparison helps
establish what changed. A source inspection supports a citation's locator
and role. Neither those checks nor a lower rate of phrases flagged by a style
detector establishes improved comprehension. A reported style comparison must
identify the versions, prose denominator and genre, and retain regressions
alongside gains. A claim of improved comprehension would require evidence
from readers performing a specified task on identified versions, with the
comparison and its limitations recorded.

The method leaves that empirical question open. It makes a particular editorial
decision inspectable: a source suggests an explanatory choice, the local
mathematics warrants the new sentence, and the review records its fate. The
practical standard remains the reader's ability to recover the argument. A
different proof may warrant a different choice.


\section{Transfer the method across research genres}\label{sec:genre-transfer}

The common task is to make the relation between a claim and its warrant
legible. What counts as a warrant depends on the field and the sentence.
For a mathematical theorem, state the domain and hypotheses, then explain
the proof's hard implication. A systems paper may describe a design, report an
implemented mechanism or compare performance under a workload. Identify which
of these is being claimed. An empirical comparison needs the tested
implementation, conditions, comparator and observed result; a proposed design
must not be presented as an implemented one. A build or artifact check concerns
the tested objects, not unrestricted performance or usability.

Levin and Redell distinguish papers about implemented systems, proposed
systems and theoretical work, and caution that evaluation criteria vary across
those classes~\cite{levin-redell}. The SIGPLAN empirical-evaluation guidance
asks authors to state claims and limitations clearly, and treats its checklist
as an aid to judgment rather than a universal score~\cite{sigplan-evaluation}.
These are genre-specific primary sources, not proof authority for any
particular system. A nearby systems paper's opening and evaluation should also
be read at their actual passages before transferring a sentence pattern.

The repository's systems manuscript gives a concrete example. It reports nine
rejections among ten deliberately false, author-selected edits in a historical
trial. After one edit escaped, a follow-up checked the intact baseline and
that escaped edit against a repair; the other nine edits were not rerun
\cite[recorded observations]{systems}. The first report describes those ten
trials, not a general $90\%$ reliability rate. The second describes the tested
repair, not a post-repair ten-out-of-ten result. The manuscript also records
missing original logs and the absence of an independent comparison. Those
qualifications determine which conclusion the observations can support; they
are not incidental implementation details.

In an expository scientific paper, similarly identify the underlying study or
derivation, the population or model, the observed or cited result, and the
author's synthesis. If the paper contributes no new experiment, say so.
Compare each inference with the cited primary study and with nearby expositions
in that field; do not turn a reported association into causation or a model's
prediction into an observation.

Sentence structure should expose these dependencies. Put the condition before
its consequence when the condition controls the claim; name the measured
quantity and comparator in the same sentence as the reported change; keep
parallel outcomes in parallel clauses. The value of this practice is local:
it lets a reader test exactly which evidence licenses each statement. It is
not a single prose template for mathematics, systems research and every
science.

\section*{Acknowledgement}
Wouter van Doorn's feedback on first-reader legibility, recorded in the project's
writing guidance, helped sharpen the attention to private terminology, unnecessary
notation and restrictive hypotheses. This acknowledges advice on exposition;
it does not attribute mathematical verification or endorsement to him.

\section*{Availability and authorship}
The writing skill, source records, lesson history and this paper's LaTeX source
are maintained in the public repository~\cite{ledger,skill}. The accompanying
editorial return documents the new-return examples separately from the accepted
historical cases. The public writing method and its review ledger record
selected guidance, with source-reading and acceptance boundaries preserved.
The worked revisions were made across papers on eight Erd\H{o}s problems.
Some agents worked from frozen packets; others integrated changes against the
live repository. Their editorial judgements do not measure how independent
mathematical readers understood the papers. The method has not been tested
for improved reader comprehension.
Will Cook built and directed the research infrastructure and reviewed
the claims when he could. AI agents did most of the research and drafting.
Cook did not independently verify every claim.
